\section{PCB v1.3: single-sided SMD board (CR-003)}
\label{sec:cr003}

Inputs: branch \texttt{change/CR-003} at \texttt{f0c016b} with the T-109 working tree (assembly text and labels only; copper identical to \texttt{f0c016b}, T-109 \nolinkurl{compare_copper.out}). Sources:
\begin{itemize}[noitemsep]
  \item change request and decisions: \ev{changes/CR-003.md},
        \ev{management/DECISIONS.md} (ADR-023 to ADR-029);
  \item electronics analysis section~8: \ev{reports/electronics/analysis.md};
  \item handoffs: \ev{management/handoffs/electronics-engineer-T-103.md},
        \ev{management/handoffs/systems-engineer-T-104.md};
  \item review T-102: \ev{reports/review/CR-003/findings.md};
  \item re-review T-108: \ev{reports/review/CR-003/findings-T-108.md};
  \item budgets and metrics: \ev{sim/systems/budgets.out},
        \ev{management/metrics.json}.
\end{itemize}
This section describes the current board and supersedes the v1.2
board-specific content of Sections~\ref{sec:architecture} to \ref{sec:validation}
where they differ.

\paragraph{Status.} Every result in this section is a design-file check, a
calculation or a simulation. \textbf{The v1.3 board has not been etched, soldered,
powered or measured, and nothing in it is hardware-validated.} ERC, DRC and the
rules checks passing do not validate the hardware.

\subsection{Reasons for the change}

The user makes boards at home by chemical etching and does not want to drill them
(CR-003, 2026-11-09). V1.2 had 133 drilled holes (\lbl{CALCULATION}, design file: drill file,
\ev{changes/CR-003.md}); v1.3 has 4, the optional M3 mounting holes
(\lbl{CALCULATION}, design file, \texttt{v1.3\_drill\_holes\_total} in
\ev{management/metrics.json}). The conversion (ADR-023, ADR-026):
\begin{itemize}[noitemsep]
  \item one copper layer, on the component side (F.Cu); surface-mount parts only;
        no vias, no copper pour, no silkscreen (the ORIG/QT legend is etched in
        copper and the rest is on the assembly map);
  \item the Pico (A1) is soldered on its castellations over an etched bare-copper
        keep-out; the Adafruit 1982 module (A2) stands on its own header pins,
        cut short and soldered onto surface lands; no sockets, so a failed Pico or
        module has to be desoldered from the surface (accepted, ADR-023);
  \item R1--R12 become 2512 chip resistors, R13--R16 and C1 become 1206, SW1 is
        the SMD B3S-1000, RV1 moves off the board on 3 wire pads (ADR-023 D-2), and
        J2 becomes 14 solder-wire pads with hot-melt strain relief (ADR-023 D-3);
  \item JP1 is a 3-pad surface solder bridge and J1 two surface pads for the
        I\textsuperscript{2}C configuration wires; a bare MPR121 was rejected
        (0.40\,mm pitch against the 0.6\,mm clearance rule, \lbl{SOURCE}, ADR-023);
  \item the firmware and the pin map do not change (REQ-031 PASS).
\end{itemize}

\begin{figure}[H]
  \centering
  \begin{minipage}[b]{0.49\linewidth}\centering
    \includegraphics[width=\linewidth,trim=230 150 230 150,clip]{board-top-v1.3.png}\\[-2pt]
    {\small (a) top view}
  \end{minipage}\hfill
  \begin{minipage}[b]{0.49\linewidth}\centering
    \includegraphics[width=\linewidth,trim=230 90 150 260,clip]{board-3d-v1.3.png}\\[-2pt]
    {\small (b) 3D view}
  \end{minipage}
  \caption{PCB v1.3 after T-103, KiCad renders, not photos
  (\nolinkurl{hardware/fabrication/images/board-top.png} and
  \nolinkurl{board-3d.png}). The Pico USB end is at the top edge; R1--R12 and the
  J2 wire pads are along the bottom. The module itself is not modelled: only its
  header pins are drawn, raised 3.0\,mm (the cut-pin length, \lbl{ASSUMPTION}),
  standing on the two rows of lands in the centre.}
  \label{fig:board-v13}
\end{figure}

\subsection{Board geometry after the review fixes (T-103)}

{\small
\begin{longtable}{@{}>{\raggedright\arraybackslash}p{0.42\linewidth}>{\raggedright\arraybackslash}p{0.26\linewidth}>{\raggedright\arraybackslash}p{0.28\linewidth}@{}}
\caption{v1.3 board values after T-103 (from \nolinkurl{management/metrics.json}). Values read from the design file are \lbl{CALCULATION} (design file), not measurements (ADR-029 D-4).}\label{tab:v13-geometry}\\
\toprule
Quantity & Value & Label \\
\midrule\endfirsthead
\toprule Quantity & Value & Label \\ \midrule\endhead
Board outline & 100 $\times$ 70\,mm & \lbl{CALCULATION} (design file; REQ-028) \\
Pico (A1) anchor & y = 1.48\,mm & \lbl{CALCULATION} (design file) \\
Pico end inside the board edge & 0.11\,mm & \lbl{CALCULATION} \\
USB receptacle face beyond the edge & 1.19\,mm (limit 1.0\,mm) & \lbl{CALCULATION} from \lbl{SOURCE} 1.30\,mm; limit \lbl{ASSUMPTION} \\
Cutting tolerance & 1.0\,mm (0.5 transfer + 0.5 trim) & \lbl{ASSUMPTION} \\
A2 lands, electrode row / S row & 1.8 $\times$ 4.2\,mm / 1.8 $\times$ 5.515\,mm & \lbl{CALCULATION} (design file) \\
A2 lands beyond the module outline & 1.73--1.775\,mm & \lbl{CALCULATION} (design file) \\
Header strip, cut pin, module board & 2.5 / 3.0 / 1.6\,mm & \lbl{ASSUMPTION} (first-build list) \\
Module underside / board top above the copper & 5.5 / 7.1\,mm & \lbl{CALCULATION} on the \lbl{ASSUMPTION}s above \\
1206 lands, R13--R16 / C1 & 1.775 $\times$ 1.8 / 1.875 $\times$ 1.8\,mm & \lbl{CALCULATION} (design file) \\
Keep-out rule areas under the Pico & 3 & \lbl{CALCULATION} (design file) \\
Copper to the nearest Pico~2 underside feature & 2.14\,mm (11 features) & \lbl{CALCULATION} (design file) \\
Minimum copper clearance between nets & 0.680\,mm (rule 0.6\,mm) & \lbl{CALCULATION} (design file) \\
Minimum copper to board edge & 0.680\,mm (rule 0.5\,mm) & \lbl{CALCULATION} (design file) \\
JP1 pad gap & 1.87\,mm & \lbl{CALCULATION} (design file) \\
Electrode pad to GND edge track (spark gap) & 1.401\,mm, all 12 & \lbl{CALCULATION} (design file) \\
SDA track to nearest electrode pad & 3.30\,mm & \lbl{CALCULATION} (design file) \\
\bottomrule
\end{longtable}
}

The Pico sits 0.11\,mm short of flush with the edge: flush would put the pin~1 and
pin~40 (VBUS) lands 0.495\,mm from the etched outline line, below the 0.6\,mm rule
(\lbl{CALCULATION}); the manager accepted this (ADR-028 D-1). The PCB supports
only the Raspberry Pi Pico (SC0915, recommended) and the Pico~2; Pico~W and Pico~2~W
are not supported until their undersides are checked (ADR-028 D-3).

\paragraph{Checks on the design files (T-103, T-104, T-105).}
\begin{itemize}[noitemsep]
  \item ERC: ``ERC messages: 0  Errors 0  Warnings 0''; DRC: ``Found 0 DRC
        violations'', ``Found 0 unconnected pads''; ``Found 0 schematic parity
        issues'' (KiCad 10.0.4, T-103 handoff).
  \item \nolinkurl{check_smd}: ``SUMMARY: 65 checks, 65 PASS, 0 FAIL'' (T-103); the
        new checks were each shown to fail on a mutated copy.
  \item \nolinkurl{pcb_rules_check}: ``SUMMARY: 26 checks, 26 PASS, 0 FAIL''; mutation
        tests ``MUTATION SUMMARY: 85 proved, 0 not proved'' (T-104).
  \item \nolinkurl{interface_check} after the T-105 documentation fixes:
        ``SUMMARY: 121 checks, 121 PASS, 0 FAIL''
        \ev{management/handoffs/documentation-engineer/T-105/interface_check.out}.
  \item \nolinkurl{budgets}: ``SUMMARY: 35 checks, 31 PASS, 4 FAIL''; the 4 FAILs are
        A3b (VBUS, REQ-018) and the three 100\,kHz latency rows, unchanged from
        earlier revisions (T-104).
\end{itemize}

\subsection{Bill of materials and cost}
\label{sec:v13-bom}

All parts are bought from Mouser (ADR-024). Prices are \lbl{SOURCE} Mouser Portugal,
2026-11-11, EUR excluding VAT \ev{sim/systems/budgets.out}.

{\small
\begin{longtable}{@{}>{\raggedright\arraybackslash}p{0.13\linewidth}>{\raggedright\arraybackslash}p{0.46\linewidth}>{\raggedright\arraybackslash}p{0.07\linewidth}>{\raggedright\arraybackslash}p{0.30\linewidth}@{}}
\caption{v1.3 core BOM (budgets revision 10).}\label{tab:v13-bom}\\
\toprule
Ref. & Part (Mouser number) & Qty & EUR, label \\
\midrule\endfirsthead
\toprule Ref. & Part (Mouser number) & Qty & EUR, label \\ \midrule\endhead
A1 & Raspberry Pi Pico SC0915 (358-SC0915) & 1 & 3.44 \lbl{SOURCE} \\
A2 & Adafruit MPR121 breakout 1982 (485-1982) & 1 & 6.84 \lbl{SOURCE} \\
R1--R12 & YAGEO RC2512FK-073K3L, 3k3 1\,\%, 2512 (603-RC2512FK-073K3L) & 12 & 1.31 at 0.109 (10+ tier) \lbl{SOURCE} \\
R13--R16 & YAGEO RC1206FR-074K7L and RC1206FR-07220RL & 4 & 0.34 \lbl{SOURCE} \\
C1 & Samsung CL31B104KBCNNNC 100\,nF 1206 (187-CL31B104KBCNNNC) & 1 & 0.12 \lbl{SOURCE} \\
SW1 & B3S-1000 (653-B3S-1000; ARATAS, formerly Omron Components) & 1 & 0.70 \lbl{SOURCE} \\
RV1 & Bourns PTV09A-4015F-B103 (652-PTV09A-4015FB103), off-board & 1 & 0.77 \lbl{SOURCE} \\
Knob & Same Sky SPK-016A (179-SPK-016A); bore fit NOT CONFIRMED & 1 & 0.33 \lbl{SOURCE} \\
--- & Single-sided copper clad 100 $\times$ 70\,mm & 1 & 1.00--3.00 \lbl{ASSUMPTION} \\
\midrule
 & Mouser lines only & & 13.84 \lbl{CALCULATION} \\
 & \textbf{Core BOM, ex VAT, ex shipping} & & \textbf{14.84--16.84} \lbl{CALCULATION} \\
 & Core BOM with 23\,\% VAT (rate \lbl{ASSUMPTION}) & & 18.25--20.71 \lbl{CALCULATION} \\
 & Shipping & & NOT CONFIRMED \\
\bottomrule
\end{longtable}
}

J1, J2, J3 and JP1 are etched copper only and are excluded from the BOM; the knob
has no schematic symbol and is listed in the assembly guide only (ADR-028 D-4;
budgets C2 PASS, ``differences []''). REQ-032 ($\le$ 30\,EUR) is PASS on the basis
ruled by ADR-027 D-1 (core BOM excluding shipping, consumables and tools), with
shipping reported as NOT CONFIRMED.

\paragraph{Supply.} Mouser showed 13 RC2512FK-073K3L in stock for the 12 needed,
next delivery 2027-05-18 (\lbl{SOURCE}, 2026-11-11; risk R-016, accepted by
ADR-027 D-3). The named alternative is Vishay CRCW25123K30FKEG: 12 $\times$ 0.424
= 5.09\,EUR, +3.78\,EUR, core 18.62--20.62\,EUR (\lbl{CALCULATION} from
\lbl{SOURCE} 0.424\,EUR). Its datasheet has not been opened, so its body size is
NOT CONFIRMED; it may be used only after \nolinkurl{check_gap} (d) and
\nolinkurl{sim_electrode} are re-run with its values (ADR-028 D-5).

\paragraph{How the resistor was chosen.} ADR-023 chose TT Electronics HPWC2512
3k0, which could not be bought in small quantities; ADR-024 chose the Vishay CMB
0207 (one-board quote about 17.6\,EUR, \lbl{CALCULATION}), with an accepted 3.2\,mm
conductor gap; ADR-025 chose the HPWC2512-3K3JT1 (one-board quote about 27.5\,EUR,
\lbl{CALCULATION}); ADR-026, the user's final decision, chose the standard YAGEO
part and accepted that its pulse capability is not proven.

\subsection{ESD analysis of the electrode path}
\label{sec:v13-esd}

The protection target is an IEC 61000-4-2 contact discharge of $\pm$4\,kV, itself
an \lbl{ASSUMPTION} (level 2; the standard was not read). The model is a
150\,pF / 330\,$\Omega$ source into the 3.3\,k$\Omega$ series resistor
\ev{reports/electronics/analysis.md} section~8 \ev{sim/electronics/calc/calc_pulse_rated.out}.

{\small
\begin{longtable}{@{}>{\raggedright\arraybackslash}p{0.44\linewidth}>{\raggedright\arraybackslash}p{0.30\linewidth}>{\raggedright\arraybackslash}p{0.22\linewidth}@{}}
\caption{ESD numbers for R1--R12 = RC2512FK-073K3L (analysis.md section 8).}\label{tab:v13-esd}\\
\toprule
Quantity & Value & Label \\
\midrule\endfirsthead
\toprule Quantity & Value & Label \\ \midrule\endhead
Ratings: max working / overload voltage & 200 / 500\,V; no pulse, surge or ESD rating published & \lbl{SOURCE} YAGEO RC\_L V.11 \\
Voltage across R at 4\,kV / 8\,kV & 3636 / 7273\,V (7.3$\times$ the overload rating at 4\,kV) & \lbl{CALCULATION} \\
Energy into R per strike, 4\,kV / 8\,kV & 1.09 / 4.36\,mJ, $\tau$ = 0.544\,\textmu s & \lbl{CALCULATION} \\
Peak power, 4\,kV / 8\,kV & 4007 / 16028\,W & \lbl{CALCULATION} \\
Old 1.2/50\,\textmu s acceptance energy (MBE0414 182\,mJ; HPWC 231\,mJ) & 167$\times$ / 212$\times$ one 4\,kV strike & \lbl{CALCULATION} \\
Film temperature rise, 4\,kV / 8\,kV & 3.1 / 12.6\,K; 31 / 126\,K with a $\times$10 trim hot spot & \lbl{CALCULATION} on \lbl{ASSUMPTION} \\
Sensor pin current at 4\,kV & 1.10\,A (1.11\,A at 3267\,$\Omega$) vs 1.33\,A inferred limit; margin 1.20--1.22 & \lbl{CALCULATION}; limit \lbl{ASSUMPTION} \\
Charge, IEC 4\,kV vs HBM 2\,kV & 600 vs 200\,nC (3.0$\times$) & \lbl{CALCULATION} \\
Decay time, IEC through 3.3\,k$\Omega$ vs HBM & 0.54 vs 0.15\,\textmu s (3.6$\times$) & \lbl{CALCULATION} \\
Pin current at 8\,kV & 2.20\,A: FAIL & \lbl{CALCULATION} \\
Copper gap across R / termination gap & 5.00\,mm / 4.65\,mm (4.08\,mm at $\pm$0.5\,mm placement) & \lbl{CALCULATION} (design file) / \lbl{CALCULATION} \\
Gap rule at 4\,kV / 8\,kV & $\ge$ 3.5 / $\ge$ 7.5\,mm: PASS / FAIL & \lbl{CALCULATION} on \lbl{ASSUMPTION} models \\
Electrode sensitivity, worst ratio at 3267--3333\,$\Omega$ & 0.904 (limit 0.90) & \lbl{SIMULATION} \\
\bottomrule
\end{longtable}
}

\paragraph{What the numbers show and do not show.}
\begin{itemize}[noitemsep]
  \item \textbf{Pulse capability of the resistor: NOT PROVEN} at 4\,kV and 8\,kV.
        The maker publishes no pulse rating. The energy comparison (one 4\,kV strike
        is 1/167 of the energy the earlier, rated part was accepted against) is not
        a rating for this part: a thick-film resistor under a fast pulse fails at the
        local field and hot spot of the laser trim, which the datasheet does not
        cover. The user accepted this to keep the board cheap (ADR-026).
  \item The pin-current margin of 1.20--1.22 compares peak current only. The 4\,kV
        IEC pulse carries 3.0 times the charge and lasts 3.6 times longer than the
        HBM pulse the sensor rating comes from, so the margin is not a proof
        (review F3-005).
  \item The spark gap of 1.401\,mm from each electrode pad to the GND edge track is
        intended to take a strong discharge before the 5\,mm gap around the
        resistor (ADR-020); it is also unproven.
  \item Only an ESD-gun test can close REQ-023 (PARTIAL, test NOT RUN).
\end{itemize}

\paragraph{Plan B and its trigger (ADR-027 D-2, risk R-018).} If, after the
physical ESD series, any R1--R12 reads outside 3234--3366\,$\Omega$ ($\pm$2\,\%,
\lbl{ASSUMPTION} window from the 0.901 sensitivity at 3366\,$\Omega$,
\lbl{SIMULATION}), a channel stops detecting touch, or an arc crosses a resistor,
the board moves to the T-097 clamp: 12 $\times$ Bourns CDSOD323-T05LC to GND with
R1--R12 at 1\,k$\Omega$ (sensitivity 0.938, \lbl{SIMULATION}; +6.18\,EUR, core
21.02--23.02\,EUR, \lbl{CALCULATION}), after its maximum capacitance and leakage
are confirmed.

\paragraph{Board-level residual items (ADR-028 D-7).} Each rests on
\lbl{ASSUMPTION} inputs and is NOT RUN \lbl{[PHYSICAL]}:
\begin{itemize}[noitemsep]
  \item GND bounce between the module and the Pico during a strike to the spark gap:
        up to 330\,V for about 1\,ns; 1.49\,A and 0.90\,nC per I\textsuperscript{2}C
        line, against 2.67\,A and 400\,nC for the RP2040 HBM 4\,kV FT-pin equivalent
        (\lbl{CALCULATION}; budgets D1 PASS; the GP26/GP27 rating is NOT CONFIRMED,
        risk R-010). The taller module pin path of T-103 raised it from 236\,V.
  \item SDA coupling to the electrode pads: 36.6\,fF, a 2.09-count step against the
        12-count touch threshold (\lbl{CALCULATION}, analysis.md section 8.5; budgets
        D2 PASS against the 3-count proposal of REQ-038; risk R-017).
  \item USB pull: 35.6\,N on the most loaded M3 screw at a 30\,N pull
        (\lbl{CALCULATION}; budgets D3 PASS against an unsourced 100\,N
        \lbl{ASSUMPTION}; risk R-013).
\end{itemize}

\subsection{Independent review T-102 and dispositions}
\label{sec:cr003-review}

The independent reviewer reviewed \texttt{9505245} and found the board \textbf{not
ready for the first home-etched build}: F3-001 blocked, F3-002 and F3-003 had to be
fixed or scoped first \ev{reports/review/CR-003/findings.md}. The fixes are T-103
(electronics), T-104 (systems) and T-105 (documentation); the manager's rulings are
ADR-028. Status below is the result of the check named for each fix on the design
files or documents. The fixes were re-reviewed in T-108 (Section~\ref{sec:cr003-rereview}).

{\footnotesize
\begin{longtable}{@{}>{\raggedright\arraybackslash}p{0.07\linewidth}>{\raggedright\arraybackslash}p{0.07\linewidth}>{\raggedright\arraybackslash}p{0.30\linewidth}>{\raggedright\arraybackslash}p{0.40\linewidth}>{\raggedright\arraybackslash}p{0.10\linewidth}@{}}
\caption{CR-003 review findings and dispositions (ADR-028).}\label{tab:review-cr003}\\
\toprule
ID & Sev. & Finding & Disposition & Status \\
\midrule\endfirsthead
\toprule ID & Sev. & Finding & Disposition & Status \\ \midrule\endhead
F3-001 & critical & A2 lands entirely under the header strip and module outline; joints cannot be reached or inspected. & Lands lengthened outward to 1.73--1.775\,mm beyond the outline; strip pushed against the module, pins cut to 3.0\,mm, soldered from outside (T-103; check\_smd land check). & PASS (design file) \\
F3-002 & major & USB face only 0.67\,mm beyond the edge. & Pico moved; face 1.19\,mm beyond, Pico end 0.11\,mm short of flush, accepted (ADR-028 D-1); oversize cut, trim and plug-seat check in the guide. & PARTIAL \\
F3-003 & major & Pico~2 TP7 outside the keep-outs; Pico~W / Pico~2~W never checked. & TP7 keep-out added, 11 Pico~2 features clear by $\ge$ 2.14\,mm; PCB scoped to Pico and Pico~2 (ADR-028 D-3); README rows corrected in T-105. & PASS (design file, documents) \\
F3-004 & major & ESD analysis described the old HPWC part. & analysis.md section 8 on RC2512 (pulse NOT PROVEN); 7.1 and 7.2 marked superseded; guide link fixed. & PASS (document) \\
F3-005 & minor & Pin margin compares peak current only. & Charge 3.0$\times$ and decay 3.6$\times$ stated beside the margin. & PASS (document) \\
F3-006 & minor & ORIG label is left of JP1, not below. & Copper unchanged; guide (T-103) and manual (T-105) describe the real position (ADR-028 D-2). & PASS (documents) \\
F3-007 & minor & No post-solder short checks for VBUS/VSYS and BTN/RUN. & Checks added: pads 36, 39, 40 not to GND; BTN not to pin 30. & PASS (document) \\
F3-008 & minor & Pin-1 mark level with pin 2. & Mark moved level with pin 1, square drawn in pad 1. & PASS (render inspected) \\
F3-009 & minor & Records disagree (REQ-032, check counts, stock date, INTERFACES). & Requirements draft-13 to 16 and INTERFACES frozen-9, proposed in T-104 and frozen by the manager (ADR-028 D-8); ADR-026 stock date superseded (ADR-028 D-6); frozen-9 frozen. & PASS (records) \\
F3-010 & minor & J1 in the BOM without MPN; knob outside the KiCad BOM; SW1 maker. & J1, J2, J3, JP1 copper only; knob in the guide only; SW1 genuine B3S-1000 (ADR-028 D-4). & PASS (budgets C2) \\
F3-011 & info & Wire direction and strip length at J2. & Wires leave upward, strip 3--4\,mm. & PASS (document) \\
F3-012 & info & Solder at the USB shell pins can tilt the Pico. & Rocking check before tacking. & PASS (document) \\
F3-013 & info & ``module'' used for the Pico. & Wording fixed. & PASS (document) \\
F3-014 & info & 1206 lands 0.43\,mm beyond the body. & Lands lengthened; $\ge$ 1.05\,mm beyond the body. & PASS (design file) \\
F3-015 & info & Regeneration not byte-identical (UUIDs). & No action; geometry identical. & --- \\
\bottomrule
\end{longtable}
}

\subsection{Independent re-review T-108 and dispositions}
\label{sec:cr003-rereview}

A second independent reviewer, who did not see the work, re-reviewed
\texttt{f0c016b} \ev{reports/review/CR-003/findings-T-108.md}. Verdict: \textbf{ready
for the first build after fixes}; no critical finding, and no copper change
needed. The reviewer's own checks on a git archive (kicad-cli 10.0.4 via pcgate)
repeat ERC 0, DRC 0 / 0 unconnected / 0 parity and the board values of
Table~\ref{tab:v13-geometry} (smallest net gap 0.680\,mm, 0.605\,mm to the etched
outline line; A2 lands 1.730 / 1.775\,mm beyond the outline; 12 Pico~2 underside
features $\ge$ 2.14\,mm from copper; all \lbl{CALCULATION} from the design file).
Review status: PARTIAL (design-file review complete; everything physical NOT RUN
\lbl{[PHYSICAL]}). The manager accepted all ten findings (ADR-029); the electronics
fixes are T-109 \ev{management/handoffs/electronics-engineer-T-109.md}
(assembly text and labels only, copper identical, ERC 0, DRC 0,
\nolinkurl{check_smd} 65/65), and this documentation update is T-111.

{\footnotesize
\begin{longtable}{@{}>{\raggedright\arraybackslash}p{0.07\linewidth}>{\raggedright\arraybackslash}p{0.07\linewidth}>{\raggedright\arraybackslash}p{0.30\linewidth}>{\raggedright\arraybackslash}p{0.40\linewidth}>{\raggedright\arraybackslash}p{0.10\linewidth}@{}}
\caption{CR-003 re-review findings and dispositions (ADR-029).}\label{tab:review-cr003-t108}\\
\toprule
ID & Sev. & Finding & Disposition & Status \\
\midrule\endfirsthead
\toprule ID & Sev. & Finding & Disposition & Status \\ \midrule\endhead
F4-001 & major & Epoxy bead at the USB end has no board to bond to (Pico end 0.11\,mm inside the edge) and would hang where the plug moulding seats. & Beads on the two long sides over pins 1--3 and 38--40; none on the end face or beyond the edge; plug-seat re-check after curing (ADR-029 D-1). Guide, manual and root README. & PASS (documents) \\
F4-002 & minor & Pad 40 is 3.48\,mm from GND pad 38, not 0.94\,mm; pad 37 (3V3\_EN) to 38 (0.94\,mm) not checked. & Sentence corrected; ``37 to 38 must not beep'' and a troubleshooting line added. & PASS (documents) \\
F4-003 & minor & About 98\,\% of the etched ORIG lies under the module, and JP1 was bridged after the module. & JP1 before the module; pads identified by position; text says ORIG is hidden; label move deferred to the next copper change (ADR-029 D-2). & PASS (documents) \\
F4-004 & minor & S-row pins 2.54\,mm inside the module outline (electrode row 1.27\,mm); joints visible only from a low angle. & Low-angle inspection described; S pad to module-pin continuity check added. & PASS (documents) \\
F4-005 & minor & Bushless PTV09A with no mounting method. & Perfboard mounting behind a 7\,mm panel hole (\lbl{ASSUMPTION}; method not tried) (ADR-029 D-3). & PASS (documents) \\
F4-006 & minor & ``Soldered flat'' is wrong for the module. & ``The module stands on its header pins.'' & PASS (documents) \\
F4-007 & info & Design-file geometry labelled MEASUREMENT. & Relabelled \lbl{CALCULATION} (design file) in check scripts, analysis, metrics, requirements, manual and this report (ADR-029 D-4). & PASS (documents) \\
F4-008 & info & Stale ``module end y 0.63''. & Marked as T-092 history (comment and analysis only). & PASS (documents) \\
F4-009 & info & 0.8\,mm S-row stubs lie inside lands and pads. & Note in the design-rule table. & PASS (document) \\
F4-010 & info & USB overhang 1.30\,mm sourced for the Pico only. & Plug-seat check made explicit for the Pico~2. & PASS (documents) \\
\bottomrule
\end{longtable}
}

Not checked by the reviewer: the S-row order of both module versions (Adafruit
Eagle files not on disk), the Pico~2 USB overhang, tenting of the Pico and Pico~2
underside vias, Mouser stock, a bushed PTV09A variant, and anything physical.

\subsection{First-build measurement list}
\label{sec:v13-first-build}

The first build has to measure the values the design rests on. All items are NOT RUN
\lbl{[PHYSICAL]}; sources are the REQ-029 and REQ-023 evidence in
\ev{requirements/requirements.csv} and the assembly guide
\ev{hardware/README.md}. If a value falls outside its assumption,
\nolinkurl{check_smd} and \nolinkurl{pcb_rules_check} are re-run with the measured
value.

{\footnotesize
\begin{longtable}{@{}>{\raggedright\arraybackslash}p{0.06\linewidth}>{\raggedright\arraybackslash}p{0.32\linewidth}>{\raggedright\arraybackslash}p{0.27\linewidth}>{\raggedright\arraybackslash}p{0.31\linewidth}@{}}
\caption{v1.3 first-build measurements (all NOT RUN).}\label{tab:v13-first-build}\\
\toprule
\# & What & Design value (label) & How / criterion \\
\midrule\endfirsthead
\toprule \# & What & Design value (label) & How / criterion \\ \midrule\endhead
FB-1 & Height of the supplied header strip & 2.5\,mm (\lbl{ASSUMPTION}) & Caliper; re-run the checks if different \\
FB-2 & Exposed pin length after cutting, all 20 & 3.0\,mm (\lbl{ASSUMPTION}) & Caliper; all equal, module level on both rows \\
FB-3 & Module board thickness; module underside and top above the copper & 1.6\,mm (\lbl{ASSUMPTION}); 5.5 / 7.1\,mm (\lbl{CALCULATION}) & Caliper \\
FB-4 & C1 body, and copper beyond it & up to 3.4 $\times$ 1.8\,mm (\lbl{ASSUMPTION}); $\ge$ 1.0\,mm & Caliper on the fitted part \\
FB-5 & B3S body & 6.0 + 0.1\,mm (\lbl{ASSUMPTION}) & Caliper \\
FB-6 & USB receptacle face beyond the trimmed edge; Pico end inside the edge & 1.19\,mm (limit 1.0\,mm); 0.11\,mm (\lbl{CALCULATION}) & Caliper after trimming; plug fully seated before soldering, explicitly with a Pico~2 (overhang sourced for the Pico only, F4-010), and again after the epoxy beads cure (F4-001) \\
FB-7 & One RC2512 body L and termination I1 & L 6.35 $\pm$0.10, I1 0.60 $\pm$0.20\,mm (\lbl{SOURCE}) & Caliper; inside 6.25--6.45 and 0.40--0.80\,mm, else re-run \nolinkurl{check_gap} (d) \\
FB-8 & R1--R12 before soldering and after the ESD series & 3267--3333\,$\Omega$ (3.3\,k$\Omega$ $\pm$1\,\%, \lbl{CALCULATION}) & Ohmmeter; after the series every part within 3234--3366\,$\Omega$, else plan B \\
FB-9 & Pico flat before tacking & --- & Must not rock at any corner \\
FB-10 & Post-solder continuity & --- & Pico pads 36, 39, 40 not shorted to GND; pad 37 not to pad 38 (F4-002); SW1 BTN not to pad 30; JP1 centre only to S1 (ORIG) or S7 (QT); each S pad to its own module-top pin and not to the neighbour (F4-004) \\
FB-11 & ESD series $\pm$4\,kV contact on E0--E11 and the RV1 knob & --- & REQ-023 plan steps (1)--(3); arc path, resets and resistor values recorded \\
FB-12 & SDA coupling step, GND bounce effects, USB pull & 2.09 counts; 330\,V / 1\,ns; 35.6\,N (Section~\ref{sec:v13-esd}) & Touch data during I\textsuperscript{2}C traffic; behaviour during strikes; no pad lift at a 30\,N pull \\
\bottomrule
\end{longtable}
}

\subsection{Requirement status for CR-003 (draft-16)}

Requirements draft-16 (T-112), 42 requirements: PASS 18, PARTIAL 18, FAIL 2
(REQ-018, REQ-040), NOT RUN 4 \ev{requirements/requirements.csv}. Rows
touched by CR-003:

{\small
\begin{longtable}{@{}>{\raggedright\arraybackslash}p{0.10\linewidth}>{\raggedright\arraybackslash}p{0.24\linewidth}>{\raggedright\arraybackslash}p{0.11\linewidth}>{\raggedright\arraybackslash}p{0.51\linewidth}@{}}
\caption{CR-003 requirements in draft-16.}\label{tab:v13-req}\\
\toprule
ID & Title & Status & Basis \\
\midrule\endfirsthead
\toprule ID & Title & Status & Basis \\ \midrule\endhead
REQ-023 & Electrode ESD robustness & PARTIAL & Analysis only; resistor pulse capability NOT PROVEN (ADR-026); ESD test NOT RUN \\
REQ-025 & Both module revisions & PARTIAL & JP1 and wires checked on the design file; test NOT RUN \\
REQ-027 & Single-layer SMD PCB & PASS & Design file (pcb\_rules\_check 26/26); no board etched \\
REQ-028 & Board outline & PASS & 100.000 $\times$ 70.000\,mm on the design file \\
REQ-029 & Home-etch and SMD land rules & PASS & Design file only; etching and hand soldering NOT RUN; first-build list above \\
REQ-030 & Electrode wire pads & PASS & Design file \\
REQ-031 & Pin map consistency & PASS & interface\_check, 0 pin mismatches \\
REQ-032 & Core BOM cost & PASS & 14.84--16.84\,EUR ex VAT, ex shipping (ADR-027 D-1); shipping NOT CONFIRMED \\
REQ-036 & User and assembly manual & PASS & Draft-16 (T-112): interface\_check 125/125, documentation 35/35, after T-111. Covers consistency of the manual and READMEs with the design files only; the hand-assembly steps have not been tried \\
REQ-040 & JP1 wrong-position tolerance & FAIL & Accepted risk R-007 (ADR-027 D-5) \\
\bottomrule
\end{longtable}
}

A PASS in this table means the analysis, inspection or design-file check passed.
None of these requirements has been verified on a built board.
